\section*{第 \ref{ch:b1:populations} 章——种群与种群统计学}

\begin{solution}{exo:b1:populations:1}
\omterm{def:b1:populations:population}{种群}：同一地点、同一时间中同一物种的个体，
彼此能够交配。密度：单位面积或单位体积内的个体数。
分布型：个体的空间格局——集群分布（一群牛、一片荨麻
），均匀分布（筑巢塘鹅相距仅一啄，石炭酸灌木），
随机分布（草坪上的蒲公英）。
\end{solution}

\begin{solution}{exo:b1:populations:2}
$l_x$：同批个体中活到年龄 $x$ 时的比例。$m_x$：
年龄为 $x$ 的雌性个体所产雌性后代数。$R_0 = \sum l_x m_x$：每只新生
雌性一生中所产雌性后代数。$T = \sum x l_x m_x / R_0$：母体平均年龄。
$R_0 = 1$：每只雌性恰好由一只雌性替代自身，\omterm{def:b1:populations:population}{种群}保持稳定。
\end{solution}

\begin{solution}{exo:b1:populations:3}
III 型：每 1000 个中约 15 个，\qty{1.5}{\%}。I 型：每
1000 个中约 700 个，\qty{70}{\%}。
\end{solution}

\begin{solution}{exo:b1:populations:4}
$\dd N/\dd t = rN$ 与 $\dd N/\dd t = rN(1 - N/K)$。$r$ 为
资源不受限时的人均增长率（单位时间内每个个体
的出生数减死亡数）；$K$ 为\omterm{thm:b1:populations:logistic}{环境容纳量}，即
增长停止时的\omterm{def:b1:populations:population}{种群}大小。
\end{solution}

\begin{solution}{exo:b1:populations:5}
$r = \ln(1000)/5 = \qty{1.38}{h^{-1}}$；倍增时间 $\ln 2/1.38 =
\qty{0.5}{h}$。
\end{solution}

\begin{solution}{exo:b1:populations:6}
$l_x m_x = 0, 0.5, 0.6, 0.2, 0$：$R_0 = 1.3$。$T = (1\times 0.5 +
2\times 0.6 + 3\times 0.2)/1.3 = 2.3/1.3 = 1.77$。$r = \ln 1.3/1.77 =
\qty{0.15}{}$ 每单位时间，\omterm{def:b1:populations:population}{种群}在增长。
\end{solution}

\begin{solution}{exo:b1:populations:7}
$N = 50\times 70/7 = 500$。仅余 40 条带标记的鱼：$N = 40\times
70/7 = 400$；第一次估计高出四分之一（标记
个体占\omterm{def:b1:populations:population}{种群}的比例小于假设）。
\end{solution}

\begin{solution}{exo:b1:populations:8}
$0.2\times 100\times 0.9 = 18$；$0.2\times 500\times 0.5 = 50$；
$0.2\times 900\times 0.1 = 18$ 每年。最大值 $rK/4 = 50$ 每年，出现于
$N = 500$。
\end{solution}

\begin{solution}{exo:b1:populations:9}
蒲公英：$r$ 选择（数千粒种子，一年生）。大象：$K$ 选择（每四年
产一头幼崽，寿命长，有育幼行为）。鳕鱼：$r$ 选择（数百万枚卵，无亲代照料）。
大猩猩：$K$ 选择（每四年一个婴儿，寿命数十年）。家蝇：$r$
选择（数百枚卵，十天内成熟）。
\end{solution}

\begin{solution}{exo:b1:populations:10}
$r = \ln(29/10)/2 = \qty{0.53}{h^{-1}}$；$K \approx 660$。在 \qty{8}{h} 时：
$660/(1 + 65e^{-4.24}) = 660/(1 + 0.94) = 340$，而计数为 351。
增长最快，此时数量为 $K/2 = 330$，在 $t = \ln 65/0.53 = \qty{7.9}{h}$
时达到，速率为 $\dd N/\dd t = rK/4 = 87\times 10^5$ 个\omterm{def:b1:cell-unit-of-life:cell}{细胞}每毫升每
小时。
\end{solution}

\begin{solution}{exo:b1:populations:11}
调节要求人均增长率随 $N$ 上升而下降，使得
$N$ 在扰动后趋于平衡；而一个按固定比例杀死个体的因子
只改变 $N$，不改变增长率
对 $N$ 的依赖，因此\omterm{def:b1:populations:population}{种群}之后恢复同样的增长。气候
对昆虫的限制方式是：打击既频繁又随机，
使\omterm{def:b1:populations:population}{种群}在接近任何上限前即被压低：一个
不受调节、被反复干扰压制的\omterm{def:b1:populations:population}{种群}，其
平均大小取决于恶劣季节的出现频率，而不取决于
某个\omterm{thm:b1:populations:logistic}{环境容纳量}。
\end{solution}

\begin{solution}{exo:b1:populations:12}
两个大小相同的\omterm{def:b1:populations:population}{种群}，一个底部宽广（幼龄个体多），另一个
年龄组分布均匀，二者未来截然不同：前者
将持续增长数十年，即使其 $m_x$ 现已降至更替水平，因为
众多幼体尚未繁殖（\omterm{def:b1:populations:population}{种群}惯性）；后者
则不然。$N$ 记录的是过去；年龄分布与
$m_x$ 随年龄的安排，才是方程作用的对象，而
年龄金字塔可预测下一代的规模，而 $N$
对此毫无说明。
\end{solution}

\begin{solution}{pb:b1:populations:1}
\textbf{1.} $r = \ln(29/10)/2 = \qty{0.53}{h^{-1}}$。
\textbf{2.} $\ln 2/0.53 = \qty{1.3}{h}$。
\textbf{3.} $K \approx 660\times 10^5 = \qty{6.6e7}{}$ 个\omterm{def:b1:cell-unit-of-life:cell}{细胞}每
毫升。
\textbf{4.} $N(t) = 660/(1 + 65e^{-0.53t})$：在 4 小时，$660/(1 + 7.8) =
75$（计数 71）；在 8 小时，340（351）；在 12 小时，$660/(1 + 0.112) = 594$
（595）。
\textbf{5.} 数量达到 $K/2 = 330$，时间 $t = \ln 65/0.53 = \qty{7.9}{h}$；
最大速率 $rK/4 = 0.53\times 660/4 = 87\times 10^5$ 个\omterm{def:b1:cell-unit-of-life:cell}{细胞}每
毫升每小时。
\textbf{6.} $10e^{0.53\times 12} = 10\times 578 = 5780$：
为观测值 595 的十倍。
\textbf{7.} 糖（或氧气）耗尽；
乙醇与酸类物质积累，毒害\omterm{def:b1:cell-unit-of-life:cell}{细胞}。
\textbf{8.} $q_0 = 0.40$；$q_4 = 1 - 0.34/0.40 = 0.15$：早期死亡率高，
之后死亡率较为稳定——介于 II 型与 III 型之间，与
多数野生大型哺乳动物一致。
\textbf{9.} $l_x m_x = 0, 0, 0.312, 0.368, 0.320, 0.272, 0.182, 0.080,
0.018, 0$：$R_0 = 1.55$。
\textbf{10.} $\sum x l_x m_x = 0.624 + 1.104 + 1.280 + 1.360 + 1.092 +
0.560 + 0.144 = 6.16$；$T = 6.16/1.55 = \qty{4.0}{yr}$。
\textbf{11.} $r = \ln 1.55/4.0 = \qty{0.11}{yr^{-1}}$；倍增时间为
\qty{6.3}{yr}。
\textbf{12.} $200e^{1.1} = 600$ 头鹿。
\textbf{13.} $N(10) = 600/(1 + 2e^{-1.1}) = 600/1.67 = 360$；\omterm{def:b1:populations:population}{种群}超过 500
的条件是 $(K - N_0)/N_0\,e^{-rt} = 0.2$，即 $e^{-rt} = 0.1$，$t = \ln
10/0.11 = 21$ 年。
\textbf{14.} $120\times 150/18 = 1000$ 条鲈鱼。
\textbf{15.} 100 条带标记的鱼：$100\times 150/18 = 833$；第一次估计
偏高（在湖中的带标记鱼少于假设）。
\textbf{16.} $rK/4 = 100$ 每年，在 $N = 500$ 时。
\textbf{17.} $T_2 = \ln 2/0.4 = \qty{1.7}{yr}$。
\textbf{18.} 捕获量 $hN$，取 $h = 0.2$：由 $rN(1 - N/K) = hN$ 得 $N^* =
K(1 - h/r) = 1000\times 0.5 = 500$，每年捕获 \num{100}——
低于最大可持续产量，但这一捕获比例会自动
随存量调整：存量减半，捕获量也随之减半，\omterm{def:b1:populations:population}{种群}
便会恢复。
\textbf{19.} 在 500 时：$0.4\times 500\times 0.5 = 100 < 120$：\omterm{def:b1:populations:population}{种群}
下降。在 300 时：$0.4\times 300\times 0.7 = 84 < 120$：
下降更快——一旦低于 $K/2$，固定的捕获量超过最大
产量便会使其走向灭绝。
\textbf{20.} 在 $K/2$ 处盈余最大，但任何误差——
例如\omterm{def:b1:populations:population}{种群}数量高估了五分之一，或一个坏年景使 $r$ 下降——
都会使可持续的捕获变为下降，而低于 $K/2$ 时盈余
随\omterm{def:b1:populations:population}{种群}缩小：该平衡点对过度
捕捞不稳定。略高于 $K/2$ 的捕捞则留有余地。
\textbf{21.} $R_0$ 减半至 $0.78 < 1$：鹿群缩小，直到
降到 400 以下，$m_x$ 恢复——一个密度制约因子，
将鹿群调节在 400 附近。
\textbf{22.} 否：它移除 \qty{40}{\%}，无论\omterm{def:b1:populations:population}{种群}大小如何，并使
人均增长率不变；此后数年，鹿群
以相同的 $r$ 从较低起点重新增长——这是干扰，而非
调节。
\textbf{23.} $0.11N(1 - N/600) = 30$：$N^2 - 600N + 163\,600 = 0$，
$N = 300\pm\sqrt{90\,000 - 163\,600}$——无实数解：每年 30 只
超过最大盈余 $rK/4 = 16.5$，狼群将
鹿群推向灭绝。若每年 15 只：$N = 300\pm\sqrt{90\,000 -
81\,800} = 300\pm 91$，即 209 与 391；较大的解稳定
（在其上方，盈余小于 15，鹿群回落；在其下方，
盈余更多，鹿群回升），较小的解不稳定。
\textbf{24.} 猞猁数量对食物的响应存在延迟：野兔数量上升
会在随后一两年提高猞猁的出生数与\omterm{def:b1:populations:lifetable}{存活率}；野兔数量骤减后，
只有在猞猁的\omterm{def:b1:lipids:triglyceride}{脂肪}储备与当季
幼崽耗尽之后，才会使猞猁挨饿。
\textbf{25.} 酵母：$K = \qty{6.6e7}{}$ 个\omterm{def:b1:cell-unit-of-life:cell}{细胞}每毫升，$r =
\qty{0.53}{h^{-1}}$。鹿：$R_0 = 1.55$，$T = \qty{4.0}{yr}$，$r =
\qty{0.11}{yr^{-1}}$。
\end{solution}
